\section{The Force handler}
\label{sec:force}
Protocol~5 is the only cross-region Force transition. It executes
explicit L1 instructions. Protocol~3 consumes its verified closure
and retires local programs separately. Silence, a health report or a
local timeout does not execute this transition.

\subsection{Explicit control and public responsibility}
\label{sec:force-stages}
An MSIG binds the chain, target branch version $\beta_{\rm target}$,
activation threshold $\theta_a$ and unique L1 control sequence $c$.
It contains no retained height and no implicit status change.
The committee submits it under \cref{ass:committee,ass:monitor}.
L1 finalization makes it known; its effect begins at $J_{q(\theta_a)}$.
Until that position, regular handoffs and valid old-version output
may continue. A regular handoff changes $r$, not $\beta$.

The \defterm{export session} is responsible for the next DATA after
$G$. Its source context, immutable quota and accepted count determine
the Force retry slot. It has no cross-region failure timer.
Its program can be Waiting while this public responsibility exists,
or a speculative later program can already be producing.

STATUS changes future eligibility only. A Force without OFF can select
the same physical region when it is the only eligible candidate.
If no region is eligible, the new boundary retains its anchor and target
slot without opening a regional engine. Later ON can enable selection.
Another valid MSIG can supersede the unassigned boundary.

\begin{protocol}{alg:force}{Apply L1 controls and execute an explicit Force}
Before the payload of authenticated finalized block $J_q$, let
$\theta=\theta_q$ and capture the entry version $\beta_0$.
\begin{enumerate}[label=\textbf{\arabic*)},itemsep=3pt]
\item \textbf{Apply independent status.}
From the complete applicable L1 control prefix, apply every due
independent STATUS not yet processed, in $(\theta_a,c)$ order.
Persist each effect. Repeated source delivery or equal timestamps
do not apply a record twice.

\item \textbf{Check due MSIGs.}
Inspect previously unprocessed due MSIGs in the same order.
Only the first valid instruction whose target equals $\beta_0$
executes. Mark every other due instruction processed without a Force.
Do not use a version created in this prelude as a fresh matching target.
If none matches, leave the branch version open and continue to step~4.
There is no output, Ready, elapsed-time or health-condition check.

\item \textbf{Close and retain.}
Capture the assigned export session's source context, index and
accepted count, or the stored source context and target slot if
unassigned. Retain all pre-payload $G$ with its original provenance.
Record the MSIG cause, triggering header, retained anchor and selection
lower bound $q$. Close $\beta_0$ and set $\beta=\beta_0+1$ exactly once.
Use \cref{sec:retry-slot} to fix the target slot. Archive the closed
version's unconfirmed suffix and regular successors.
The new engine, when entered, starts with fresh core state and $g=t=0$.
No new-version DATA is accepted in $J_q$ itself.

\item \textbf{Snapshot and resolve.}
Save the pre-payload eligibility snapshot $\chi_E(q)$ after all
control effects. For a new or pending boundary, resolve its first
eligible historical NextLive position under the designated authority.
Missing required evidence pauses at this canonical operation; it does
not prove an empty candidate set. With complete evidence but no
candidate, retain the unassigned boundary and continue source replay.
When selection succeeds, fix the handoff and open its export session
with accepted count zero. No delivery deadline is created.
Then return to ordered payload processing.
\end{enumerate}
\end{protocol}

The cutoff and retained anchor are fixed even if target evidence is
obtained later. Resume the original canonical operation; do not rerun
an already persisted status change or Force. A control can have a fixed
effect while successor selection still waits for evidence.

Every prelude executes at most one Force. This does not impose a gap
between Forces in different blocks, including blocks with equal
timestamps. Each instruction must still satisfy its own authorization,
advance-announcement and target-version conditions.

\subsection{Status, stale controls and the cutoff}
An independent OFF beside a stale MSIG still takes effect. It does not
revoke a handoff already fixed at an earlier eligible position.
The stale MSIG does not close the current version or select a new anchor.
The controller must account for these separate meanings when scheduling
its operations.

All control effects precede the triggering block's payload. DATA
accepted in earlier blocks remains in $G$. An old-version result in the
triggering block cannot cancel a Force by providing late progress.
It may be retained as regional evidence but cannot append that suffix.
A newer version starts only from the fixed retained anchor.

A proof of closure contains the actual MSIG and complete applicable
source context. A bare timestamp, an unavailable RPC response, an
unfinalized control or a monitoring alarm is not such a proof.
After learning closure, regional managers apply \cref{sec:retirement};
the cutoff does not wait for their local cleanup.
